% Value-at-Risk Three Ways — research-note PDF.
%
% PROSE AND CODE ARE VERBATIM from the website article:
%   site/src/content/articles/var-three-ways.tsx
% This article interpolates its computed results directly into the prose
% (nObs, VaR levels, breach counts, Kupiec statistics). Every interpolated
% value here is the one the site renders, read from
% data/var-three-ways.ts — e.g. {d.params.nObs} -> 3,805, pc(v99.histVar) ->
% 3.42%. Code cells are inline in the article rather than in a -code.ts module,
% and carry the same interpolations in their comments.
%
% CHARTS generated by figures/make_var_three_ways.py from that same module.
\documentclass{pyportfolios-note}

\noteshorttitle{Value-at-Risk Three Ways}

\begin{document}
\thispagestyle{notecover}

\notemasthead{Quantitative Insights}{Tutorial}{Q3 2026}{}

\notetitle{Value-at-Risk,\\Three Ways}

\notedeck{Three standard recipes for the same number — and a backtest that
rejects all of them.}

\notelede{How much could we lose tomorrow? Every trading desk has to answer that
question before the market opens, and Value-at-Risk is the industry's standard
answer: one number that says ``on 99 days out of 100, the loss stays below this
line.'' The catch is that there are three standard recipes for computing that
line — and they disagree exactly when it matters, on the bad days. We compute
all three on fifteen years of real DAX data, put every one of them on trial out
of sample, and then recompute the historical quantile with Polars and DuckDB to
show what the modern data stack changes (answer: the scaling, never the number).}

\notecard
  {Risk Management}
  {NumPy · SciPy · Pandas · Polars · DuckDB}
  {DAX (\^{}GDAXI)}
  {Jan 2010 – Dec 2024}
  {Setting and defending a daily loss limit — and proving it holds up}

\begin{multicols}{2}
\begin{notepipeline}
\item Download DAX (\^{}GDAXI) closes, 2010–2024 (yfinance)
\item Historical simulation: empirical 1\% and 5\% quantiles
\item Parametric: normal, then Student-t with fitted df
\item Monte Carlo: 200,000 seeded draws from the fitted t
\item Backtest all four with a rolling 250-day window + Kupiec POF
\item Re-do the quantile in Polars expressions and DuckDB SQL
\end{notepipeline}
\end{multicols}

% ============================================ 1 The number, and the data =====
\noteneedroom{170bp}
\notesection{1.\notenumsep The number, and the data}

\begin{multicols}{2}
Line up fifteen years of daily returns from worst to best. \textbf{VaR} at 99\%
is a marker planted one percent of the way in from the bad end: the loss you do
not expect to exceed on 99 days out of 100. Formally it is the (1−α)-quantile of
returns, negated — but ``the flood line'' is the right mental picture.
\textbf{CVaR} (expected shortfall) answers the question a good risk manager asks
next: \textbf{and when the water does cross the line, how deep does it get?} It
is the average loss given a breach, and we report it alongside every method
because it is the number you actually want for limits. Our laboratory is the DAX
from 2010-01-04 to 2024-12-30: 3,805 daily returns spanning the euro crisis,
COVID and the 2022 energy shock. The worst day in the sample is 2020-03-12 at
\texttt{-12.24\%} — keep that number in mind while the normal distribution tells
you it is impossible.
\end{multicols}

\notecallout{Definitions}{VaRα = −q₁₋α(returns)  ·  CVaRα = −E[ r | r ≤ q₁₋α ] —
both quoted as positive percentages of portfolio value, one-day horizon
throughout.}

% ================================================= 2 Historical simulation ===
\noteneedroom{170bp}
\notesection{2.\notenumsep Historical simulation}

The first recipe refuses to assume anything. Sort the returns, walk one percent
of the way in from the worst end, read off the number — that is the entire
method. On the full sample the 99\% VaR is \texttt{3.42\%} and the 99\% CVaR —
the average of the worst 38 days — is \texttt{4.60\%}. At 95\% the pair is
1.94\% / 2.93\%.

\begin{notecode}
import numpy as np

def hist_var_cvar(x, alpha=0.99):
    q = np.quantile(x, 1 - alpha)       # 1% quantile of returns
    return -q, -x[x <= q].mean()        # VaR, CVaR

hist_var_cvar(ret.values)               # (0.034169, 0.045981)
\end{notecode}

That honesty cuts both ways. The method treats a sleepy 2017 Tuesday exactly
like March 2020, and it has no imagination: it cannot warn you about any loss
larger than one it has already lived through. History is its only teacher — a
problem, because the worst day of the next fifteen years is under no obligation
to have a precedent in the last fifteen.

% ================================ 3 Parametric: normal, then Student-t =======
\noteneedroom{190bp}
\notesection{3.\notenumsep Parametric: normal, then Student-t}

The second recipe trades honesty for a formula: assume returns follow a known
distribution, and VaR drops straight out of its quantile function. Assume a
normal ($\mu = 0.039\%$, $\sigma = 1.228\%$ daily) and the 99\% VaR comes out at
\texttt{2.82\%} — roughly 60bp \textbf{below} the empirical quantile. That
shortfall is the price of the assumption: the bell curve simply does not believe
in days as bad as the ones the DAX has actually had. So keep the formula but
change the bell. Refit with a Student-t — the normal's heavy-tailed cousin — and
maximum likelihood hands you df = \texttt{3.22}: violently non-Gaussian tails.
The t's 99\% VaR of \texttt{3.40\%} lands almost exactly on the historical
number, and its CVaR (5.13\%) is fatter still. (One desk convention worth
knowing: at the one-day horizon many shops zero out μ entirely — at 0.039\% a
day it is noise against σ, and estimating it adds error without information. We
keep it for completeness.)

\begin{notecode}
from scipy import stats

mu, sd = ret.mean(), ret.std(ddof=1)
z = stats.norm.ppf(0.01)
var_normal = -(mu + sd * z)                     # 2.82%

df, loc, scale = stats.t.fit(ret.values)        # df = 3.22
var_t = -(loc + scale * stats.t.ppf(0.01, df))  # 3.40%
\end{notecode}

\notefigure{figures/var-three-ways/fig1-histogram.png}
  {3,805 DAX daily returns — fitted Student-t overlay, 99\% VaR markers}
  {The t and historical markers nearly coincide; the normal sits well inside
   them}
  {Source: Author calculations, from the article's computed results. DAX daily
   returns, January 2010 – December 2024.}

Where did the normal lose those 60 basis points? Watch it happen. Below are the
three densities over the loss region from −7.5\% to −1.5\%: the normal (amber)
runs out of probability almost immediately — by −3\% it has essentially declared
such days impossible — while the fitted t tracks the kernel estimate of the real
data the whole way down. The 60bp gap in VaR is this picture, integrated.

\notefigure{figures/var-three-ways/fig2-tail.png}
  {Left-tail densities, −7.5\% to −1.5\% — normal vs Student-t vs empirical (KDE)}
  {The normal density collapses toward zero beyond −3\% while the t stays close
   to the empirical estimate}
  {Source: Author calculations, from the article's computed results.}

% =========================================================== 4 Monte Carlo ===
\noteneedroom{170bp}
\notesection{4.\notenumsep Monte Carlo}

The third recipe replaces formulas with brute force: simulate 200,000 one-day
scenarios from the fitted t (seed 42) and read the tail of the simulation as if
it were history you simply have not lived yet. On a single linear asset this is a
correctness check more than a method — it must reproduce the analytic t to
Monte-Carlo error, and it does: \texttt{3.39\%} vs the analytic 3.40\%. So why
keep it? Because simulation does not care whether a closed form exists. The
moment the book contains options, path dependence, or anything else with no
analytic quantile, it is the only recipe still standing.

\begin{notecode}
rng = np.random.default_rng(42)
draws = loc + scale * rng.standard_t(df, 200_000)

var_mc  = -np.quantile(draws, 0.01)             # 3.39%
cvar_mc = -draws[draws <= np.quantile(draws, 0.01)].mean()   # 5.08%
\end{notecode}

% =================================================== 5 The backtest decides ==
\noteneedroom{190bp}
\notesection{5.\notenumsep The backtest decides}

Three recipes, three different answers — so who is right? Here is the beautiful
thing about VaR: it is a \textbf{falsifiable forecast}. If a 99\% VaR is honest,
tomorrow's loss should exceed it on about 1\% of days — no more, no fewer — and
we can simply count. We re-estimate each method on a rolling 250-day window (t
and Monte Carlo refit every 20 days, desk-style), forecast one day ahead — 3,555
out-of-sample forecasts per method — and test the breach count with Kupiec's
proportion-of-failures likelihood ratio.

\notefigure{figures/var-three-ways/fig3-rolling.png}
  {DAX daily returns vs the rolling 250d historical 99\% VaR line, 2010-12-23 → 2024-12-30}
  {The threshold plunges after 2011, 2020 and 2022 as crash days enter the
   estimation window}
  {Source: Author calculations, from the article's computed results.}

The tell-tale shape: the VaR line jumps \textbf{after} each crisis enters the
window and relaxes as it leaves. That lag is why breaches cluster — the first
breach of the backtest lands on 2011-03-15, deep in the euro crisis, and the
worst runs come in March 2020 before the window has learned what COVID is.

\begin{notetable}{@{}p{130bp}>{\raggedleft\arraybackslash}p{62bp}>{\raggedleft\arraybackslash}p{66bp}>{\raggedleft\arraybackslash}p{92bp}>{\raggedleft\arraybackslash}p{62bp}>{\raggedleft\arraybackslash}p{60bp}@{}}
\thcell{Method} & \thcell{99\% VaR} & \thcell{99\% CVaR} & \thcell{Breaches (exp. 35.6)} & \thcell{Kupiec LR} & \thcell{p-value} \\
\theadrule
Historical & 3.42\% & 4.60\% & 61 & 15.16 & 9.9e-5 \\ \rowrule
Parametric normal & 2.82\% & 3.24\% & 81 & 43.10 & 5.2e-11 \\ \rowrule
Parametric t (df 3.22) & 3.40\% & 5.13\% & 59 & 13.04 & 0.0003 \\ \rowrule
Monte Carlo (t) & 3.39\% & 5.08\% & 59 & 13.04 & 0.0003 \\
\end{notetable}

\notefigure{figures/var-three-ways/fig4-breaches.png}
  {Out-of-sample 99\% VaR breaches per method, 3,555 forecasts — dashed line = expected count}
  {Every method breaches more often than its own 1\% promise; the normal by the
   widest margin}
  {Source: Author calculations, from the article's computed results.}

Read it honestly: \textbf{every} bar clears the expected line — even the best
method understated its own failure rate. The t roughly halves the normal's
excess, but Kupiec rejects all four at the 1\% level. And Kupiec is the lenient
examiner: it only counts breaches, never asking \textbf{when} they arrived. One
look at the chart shows them arriving in volatility clusters, which is exactly
the pattern Christoffersen's (1998) independence test is built to punish and the
pattern behind Basel's traffic-light backtest zones. The diagnosis is structural,
not a fixable bug: a rolling window only ever looks backwards, so it is late to
every regime change by construction. That is not a reason to despair; it is the
empirical case for conditional risk models — a GARCH filter rescales the tail to
\textbf{today's} volatility, and filtered historical simulation is the desk
standard for precisely this failure mode.

% ============================ 6 The same quantile in Polars and DuckDB =======
\noteneedroom{190bp}
\notesection{6.\notenumsep The same quantile in Polars and DuckDB}

Strip the finance away and historical VaR is one line of data engineering: a
quantile over a column. That makes it a perfect specimen for a question every
quant team eventually asks — does the modern data stack change the answer?
\textbf{Polars} evaluates a lazy expression pipeline over the CSV;
\textbf{DuckDB} runs SQL straight against the file. Neither needs the data in
pandas, and both stream — the identical two lines still work when ``one index,
fifteen years'' becomes ``every book in the firm, tick by tick''.

\begin{notecode}
import polars as pl
import duckdb

q_pl = (pl.scan_csv("dax_returns.csv")
          .select(pl.col("ret").quantile(0.01, interpolation="linear"))
          .collect()
          .item())                               # -0.034169015644486

q_db = duckdb.sql(
    "SELECT quantile_cont(ret, 0.01) FROM read_csv('dax_returns.csv')"
).fetchone()[0]                                  # -0.034169015644486
\end{notecode}

\begin{notetable}{@{}p{170bp}>{\raggedleft\arraybackslash}p{200bp}>{\raggedleft\arraybackslash}p{110bp}@{}}
\thcell{Engine} & \thcell{1\% quantile of returns} & \thcell{|diff| vs pandas} \\
\theadrule
pandas / NumPy & -0.034169015644486 & — \\ \rowrule
Polars 1.43 (lazy) & -0.034169015644486 & 0.0e+0 \\ \rowrule
DuckDB 1.5 & -0.034169015644486 & 0.0e+0 \\
\end{notetable}

All three agree to the last printed digit (max absolute difference
\texttt{0.0e+0}, well inside the 1e−12 tolerance the pipeline asserts). On 3,805
rows the timings are dominated by engine start-up — NumPy answers in a fraction
of a millisecond, Polars in a few, DuckDB in tens — so at this size the engines
buy you nothing but ergonomics. The crossover comes when the file stops fitting
in memory: the pandas version dies, the Polars and DuckDB versions do not change
by a character.

\notecallout{Practitioner take}{A Student-t beats a normal for VaR almost always
— here it is the difference between 81 and 59 breaches on the same data. Prefer
CVaR over VaR for limits: it sees how bad the breach is, and at 99\% it runs more
than a full point above VaR on the DAX (3.42\% → 4.60\% historical, 3.40\% →
5.13\% under the t). And never ship a VaR you have not backtested — every method
in this article looked respectable until Kupiec counted its breaches.}

\noteneedroom{200bp}
\begin{notereferences}
\item Jorion, P. \textit{Value at Risk: The New Benchmark for Managing Financial
      Risk}. McGraw-Hill.
\item Kupiec, P. (1995). Techniques for Verifying the Accuracy of Risk
      Measurement Models. \textit{Journal of Derivatives}, 3(2).
\item Christoffersen, P. (1998). Evaluating Interval Forecasts.
      \textit{International Economic Review}, 39(4), 841–862.
\item McNeil, A., Frey, R. \& Embrechts, P. \textit{Quantitative Risk
      Management: Concepts, Techniques and Tools}. Princeton University Press.
\item Companion notebook: \texttt{var-three-ways.ipynb} — reproduces every
      number from raw data (seed 42), including the Polars and DuckDB cells.
\end{notereferences}

\end{document}
